\booksubheading{Function Contract}

একটি Function তৈরি করার সময় আগে নির্ধারণ করা হয়—

\begin{center}

\textbf{কী Input নেবে}
\(\longrightarrow\)
\textbf{কী Processing করবে}
\(\longrightarrow\)
\textbf{কী Output দেবে}

\end{center}

\begin{definitionbox}

\textbf{Function Contract:}
কোনো Function \textbf{কী Input গ্রহণ করবে, কী কাজ বা Processing করবে
এবং কী Output/Return Value প্রদান করবে}—এ সম্পর্কে নির্ধারিত নিয়ম বা
শর্তকে \textbf{Function Contract} বলা হয়।

\end{definitionbox}

%================================================
% Contract Example
%================================================

\begin{codeoutputbox}{Function Contract-এর উদাহরণ}

\noindent
\begin{minipage}[t]{0.53\linewidth}
\codeheading

\begin{english}
\ttfamily
def square(n):\\
\hspace*{4mm}return n * n\\[0.5mm]
print(square(5))
\end{english}

\end{minipage}
\hfill
\begin{minipage}[t]{0.43\linewidth}
\outputheading

\begin{center}
\textenglish{\ttfamily 25}
\end{center}

এখানে—

\textenglish{\ttfamily n}
\(\longrightarrow\)
Input

\medskip

\textenglish{\ttfamily n * n}
\(\longrightarrow\)
Processing

\medskip

Return করা Square
\(\longrightarrow\)
Output

\end{minipage}

\end{codeoutputbox}

\begin{tabularx}
{\linewidth}
{
|>{\bfseries\RaggedRight\arraybackslash}p{0.29\linewidth}
|Y|
}

\hline
\rowcolor{softBlue}

\textbf{বিষয়}
&
\textbf{তথ্য}
\\
\hline

Function
&
\textenglish{\ttfamily square()}
\\
\hline

Input
&
একটি Number
\textenglish{\ttfamily n}
\\
\hline

Processing
&
\textenglish{\ttfamily n * n}
\\
\hline

Output
&
Number-এর Square
\\
\hline

\end{tabularx}

\begin{rememberbox}

\begin{center}

\textbf{Function Contract}
=
\textbf{Input + Processing + Expected Output}

\end{center}

\end{rememberbox}

%================================================
% Contract and IPO
%================================================

\booksubheading{Function Contract ও IPO Model-এর সম্পর্ক}

Function Contract-কে IPO Model-এর সঙ্গেও মিলিয়ে বোঝা যায়।

\begin{tabularx}
{\linewidth}
{
|>{\bfseries\centering\arraybackslash}p{0.28\linewidth}
|Y|
}

\hline
\rowcolor{softBlue}

\textbf{IPO অংশ}
&
\textbf{Function-এর অংশ}
\\
\hline

Input
&
Parameter / Argument
\\
\hline

Processing
&
Function Body-এর Calculation বা Logic
\\
\hline

Output
&
Return Value
\\
\hline

\end{tabularx}

\begin{codeoutputbox}{IPO ভিত্তিক Function}

\noindent
\begin{minipage}[t]{0.54\linewidth}
\codeheading

\begin{english}
\ttfamily
def area(length, width):\\
\hspace*{4mm}return length * width\\[0.5mm]
print(area(5, 4))
\end{english}

\end{minipage}
\hfill
\begin{minipage}[t]{0.42\linewidth}
\outputheading

\begin{center}
\textenglish{\ttfamily 20}
\end{center}

\textenglish{\ttfamily 5, 4}
\(\longrightarrow\)
Input

\medskip

\textenglish{\ttfamily length * width}
\(\longrightarrow\)
Processing

\medskip

\textenglish{\ttfamily 20}
\(\longrightarrow\)
Output

\end{minipage}

\end{codeoutputbox}

%================================================
% Why contract
%================================================

\booksubheading{Function Contract কেন গুরুত্বপূর্ণ?}

\begin{compacttwocolpoints}

\item Function কী কাজ করবে তা পরিষ্কার হয়।

\item কী ধরনের Input দিতে হবে বোঝা যায়।

\item কী Result আশা করা যায় তা জানা যায়।

\item Function আলাদাভাবে Test করা সহজ হয়।

\item বড় Problem-কে ছোট Function-এ ভাগ করা সহজ হয়।

\end{compacttwocolpoints}

%================================================
% Valid input
%================================================

\booksubheading{Valid Input-এর বিষয়}

Function Design-এর সময় কোন Input গ্রহণযোগ্য তা-ও বিবেচনা করা প্রয়োজন।

\begin{codeoutputbox}{Input-এর শর্ত}

\noindent
\begin{minipage}[t]{0.52\linewidth}
\codeheading

\begin{english}
\ttfamily
def divide(a, b):\\
\hspace*{4mm}return a / b\\[0.5mm]
print(divide(10, 0))
\end{english}

\end{minipage}
\hfill
\begin{minipage}[t]{0.44\linewidth}
\outputheading

এখানে Error হবে।

কারণ Division-এর ক্ষেত্রে Divisor
\textenglish{\ttfamily 0}
হতে পারে না।

\medskip

অর্থাৎ Function-এর Contract-এর অংশ হিসেবে
Valid Input-এর শর্তও বিবেচনা করা দরকার।

\end{minipage}

\end{codeoutputbox}

%================================================
% Scope
%================================================

